\subsubsection{智能体协作过程和技能包}
\label{sec:llm}

\paragraph{协作方式。}开发阶段我们使用 OpenAI Codex 编码智能体，作品视频、海报与技术报告的整理使用 Claude Code 与 Codex。队员负责提出需求和约束、接线与上下电、万用表与示波器测量、摆放被测物，以及确认屏幕、报警灯、打印纸和喇叭的实际输出；智能体负责阅读源码、原理图、实现报告和运行日志，提出可以被证伪的假设，修改代码、执行构建与测试，并整理版本和证据。是否采纳某个方案、某项功能是否算完成，由队员结合实物状态决定。

\paragraph{智能体工作流设计。}经过多轮迭代，协作固定为图~\ref{fig:agentflow} 的流程。每个节点规定输入、执行方式、产物和放行条件，失败时回到指定节点。分工上，模型负责分流和解释，脚本负责确定性的比对和检查，队员负责上板和实物结论；涉及接线和供电的动作只由队员完成，智能体不把远程命令的返回值当作物理动作已经发生。每轮结束时记录版本、命令、退出码、产物路径、验证结果和未决事项，模型提出的修改与实际采用的步骤分开登记，以保留纠错轨迹。

\begin{figure}[htbp]
\centering
\resizebox{\linewidth}{!}{%
\begin{tikzpicture}[x=1cm,y=1cm]
\node[blk=psc,minimum width=2.7cm] (n1) at (0,0) {定义任务\\验收条件、预算、基线};
\node[blk=psc,minimum width=2.7cm] (n2) at (3.4,0) {功能验证\\参考实现 vs 设备};
\node[blk=psc,minimum width=2.7cm] (n3) at (6.8,0) {实现验证\\资源、时序、约束};
\node[blk=psc,minimum width=2.7cm] (n4) at (10.2,0) {实物验收\\采样、外设、装配};
\node[blk=psc,minimum width=2.7cm] (n5) at (13.6,0) {交付\\清单与哈希};
\node[blk=espc,minimum width=2.7cm] (m) at (3.4,-1.9) {模型\\分流、解释、假设};
\node[blk=ownc,minimum width=2.7cm] (s) at (6.8,-1.9) {脚本\\比对、审计、复跑};
\node[blk=extc,minimum width=2.7cm] (h) at (10.2,-1.9) {队员\\接线、上板、实物确认};
\draw[arr] (n1)--(n2); \draw[arr] (n2)--(n3); \draw[arr] (n3)--(n4); \draw[arr] (n4)--(n5);
\draw[arr,dashed] (n2.north) -- ++(0,0.45) -| node[pos=0.25,above,lbl]{差异按键定位、修正后同输入重跑} (n1.north);
\draw[arr,dashed] (n4.north) -- ++(0,0.75) -| node[pos=0.25,above,lbl]{软件通过但实物异常：按接口层、传输层、输出层分流} (n2.north east);
\draw[thick,draw=mute] (m.north) -- (n2.south);
\draw[thick,draw=mute] (s.north) -- (n3.south);
\draw[thick,draw=mute] (h.north) -- (n4.south);
\end{tikzpicture}}
\caption{智能体协作工作流}
\label{fig:agentflow}
\end{figure}

\paragraph{协作记录。}表~\ref{tab:cases} 汇总了开发中有代表性的协作片段。提示词为队员原话，模型回答中加引号的为原话，其余为过程摘要；每个案例都给出推翻原判断的证据和修正后的结论。

\begin{table}[htbp]
\centering
\caption{代表性协作片段}
\label{tab:cases}
\footnotesize
\renewcommand{\arraystretch}{1.3}
\begin{tabularx}{\linewidth}{@{}L{1.4cm}L{3.5cm}YL{3.1cm}@{}}
\toprule
\tabhead \tabh{主题} & \tabh{队员提示词} & \tabh{模型回答与纠错} & \tabh{结论与沉淀} \\
\midrule
资源划分 & “我们不打算用深度主负责测量了，双目深度只是一个辅助功能” & 依据实现报告精简独立 PL 深度核，保留双路采集、三路 VDMA、显示、PDM 和灯光链路，深度迁到 PS；资源按功能阶段分别列出，不当作同条件优化 & LUT 13655→8574，BRAM 41.5→30.5；沉淀为分工评审与优化对照预检 \\
频率搜索 & “你把搜索上限直接放宽到30HZ”；“不要动摄像头链路” & 先指出 45 帧/s 无法区分 30 Hz 与混叠；按约束只移除约 20.3 Hz 的旧截断，合成信号回归 7、21.6、22 Hz & 30 Hz 是搜索参数，采样边界仍约 22.5 Hz；沉淀为采样检查脚本 \\
打印压扁 & “现在的Nano是新的固件，你不要改”；“有的字体被压缩了” & 保留 Nano 固件，为打印新增独立 UART；先确认渲染位图比例正常，再改为 24 行分段、排空并间隔 120 ms & 同一报告重打正常；沉淀为外设分层诊断清单 \\
语音内容 & “你刚才的试播是噪音” & “刚才只确认了播放命令被接受，没有验证声音内容。” 撤回“播放成功”的结论，转查音频格式与文件 & 区分命令成功与内容正确两类证据 \\
音量控制 & “数字会变，声音不变” & 用同一条正常语音做 0\% 与 40\% 对照，队员确认“第一遍无声、第二遍有声” & 音量传递有效，作用于下一条播放；界面增加试听入口 \\
双目近距 & “右边的图有大量无效测量点”；“这张图片是调整之前的，你重新接一张” & 询问目标距离（“不到20厘米”），旧截图被队员否定后按 20 cm、50 cm 重新取证 & 按工作距离分别保存图像和匹配结果 \\
报警链路 & “能不能解除调整测量框之后就没办法标定的这个限制” & 部署后提示“当前测量框没有对应的停机基线”；“程序已输出红灯和蜂鸣器控制信号。实物红灯现在是否亮起？” 队员：“红灯亮，蜂鸣器响” & 标定比例与停机基线是两类有效性条件；逻辑输出与实物观察分开记录 \\
时序复核 & “按照最新工程，这些数据应该是多少” & 读取最新布线报告：WNS +0.756 ns，同时列有 592 个无时钟项和 1132 个未约束端点 & 表述限定为“已约束路径时序通过”；沉淀为 Vivado 报告审计脚本 \\
视频时长 & “然后把控一下时长，控制在3分钟。” & 按场景计划渲染为 179.6 s；交付前用 ffprobe 检查发现容器时长 181.6 s，原因是旁白音轨比画面长；音轨按总时长截断后重新封装 & 交付物以实测参数验收，不以渲染计划为准 \\
\bottomrule
\end{tabularx}
\end{table}

频率搜索问题的关键对话如下。

\begin{notebox}[频率搜索范围：对话原文]
\small
\textbf{背景：}换用约 1320 r/min 的电机后，系统多次报告测量无效。智能体读取实际记录，确认采集约 45.05 帧/s、无丢帧，保存的位移频谱在约 21.6 Hz 处有明显峰值，而原算法只搜索到 $0.45f_s\approx20.3$ Hz。\par\smallskip
\textbf{队员：}“你把搜索上限直接放宽到30HZ”\par
\textbf{智能体：}“要可靠搜索到30Hz，实际采集帧率必须高于60帧/秒；当前45帧/秒无法区分30Hz信号和混叠信号。我先核查现有相机链路能否提高帧率，再确定30Hz的实现方式。”\par
\textbf{队员：}“不要动摄像头链路，就单纯放宽那个上限，不要出现测到了对的频率又被舍去的情况”\par
\textbf{智能体：}“明白，不动摄像头链路，只把搜索上限设为30Hz，移除原来约20.3Hz的截断。这样已采到的21.6～22Hz峰不会再被旧上限排除；45帧/秒本身的采样限制仍然存在。”\par\smallskip
\textbf{验证：}以生产算法做合成位移回归，7 Hz、21.6 Hz、22 Hz 在两个初始相位下通过频率与幅值检查，丢帧和基线采样率不匹配仍被拒绝。
\end{notebox}

语音案例则是一次典型的自我纠错：智能体最初依据“播放命令被接受、程序没有报错”判断试播成功，队员指出实际是整段杂音后，智能体立即撤回结论，并把证据分成四类分别验收：命令成功、数据格式正确、实际内容正确、用户体验符合预期。此后所有外设功能的验收都按这四层记录，音量问题也是用同一条已确认正常的语音做 0\% 与 40\% 对照才得出结论。

\paragraph{技能包的提炼过程。}按照赛题要求，我们把协作中反复出现的判断方法提炼为 \code{skill/} 目录下的技能包（表~\ref{tab:skills}）。每项技能保留适用任务、输入文件、执行步骤和输出格式；将常用判断写成脚本，并用正例、反例和错误输入检查退出码。这样，后续遇到相同问题时，可以直接复用已有的检查方法。

\begin{table}[htbp]
\centering
\caption{技能包内容、来源与验证}
\label{tab:skills}
\footnotesize
\renewcommand{\arraystretch}{1.28}
\begin{tabularx}{\linewidth}{@{}L{2.7cm}L{1.8cm}L{3.5cm}Y@{}}
\toprule
\tabhead \tabh{技能} & \tabh{类型} & \tabh{来源} & \tabh{已验证效果} \\
\midrule
FPGA 大模型协作 & 提示词工作流 & 全部协作案例 & 划分、契约、纠错、PPA、验收五类提示词模板与交接单 \\
软硬件划分评审 & 提示词工作流 & PL 深度核迁移 & 按功能列出资源开销，复现 LUT 13655→8574 的划分过程 \\
测量证据复核 & 提示词工作流 & 帧率混用、纹理不足、名义标定 & 从采样率、跟踪质量、基线和标定四项检查结果有效性 \\
机电装配复核 & 提示词工作流 & 预览件与打印件混用 & 打印清单与遮挡项标记与交付记录一致 \\
定点黄金比对 & 校验脚本 & 双目核 83 次 PS/PL 逐位核对 & 合成样例分别得到通过、差异与缺项、格式错误三种退出码 \\
Vivado 报告审计 & 校验脚本 & 正 WNS 被误写为“时序全部收敛” & 从报告提取资源与 WNS/WHS，并识别 592 个无时钟项 \\
采样检查 & 校验脚本 & 22 Hz 峰被 $0.45f_s$ 截断 & 45/22/30 通过并提示近边界，45/30/30 拒绝 \\
接口、优化、复现对照 & 校验脚本 & 地址表混用、跨阶段对比、版本错配 & 三组正例退出 0，反例退出 1 并指出字段，缺字段退出 2 \\
交付校验 & 校验脚本 & 打包后文件与版本错配 & 逐文件 SHA256，拒绝路径穿越与非 ASCII 文件名 \\
外设分层诊断 & 踩坑清单 & 打印压扁、杂音、音量、触屏 & 打印分段修复、音量对照的排查顺序 \\
通用 PYNQ Overlay & 案例模板 + 校验脚本 & 本项目的寄存器契约与 DMA 调试经验 & 契约校验、板端测试脚本生成、HWH 实例预检、RTL 位精确仿真 \\
\bottomrule
\end{tabularx}
\end{table}

通用 PYNQ 技能围绕 Overlay、寄存器和简单模式 AXI DMA 组织。JSON 契约记录 DMA 与 IP 实例、数据类型、长度和控制偏移；脚本据此生成板端测试程序，并预检 HWH 中的实例名。配套示例提供确定性输入、黄金输出和带背压的 AXI4-Stream 测试平台，便于定位长度、TLAST 和数据排列问题。换用其他算子时修改 RTL 与黄金函数，换板时重新生成对应的 bit/HWH 文件。现有验证包含离线契约检查和 RTL 仿真。

\FloatBarrier
